\begin{proof}[Proof of \cref{obj:lem:compact-gaussian-convolution-injective}]
\begin{enumerate}
\item Define the finite nonnegative Borel measures
\[
\lambda_{\mathrm L}:=\mu_1+\nu_2,
\qquad
\lambda_{\mathrm R}:=\nu_1+\mu_2.
\]
Distributivity of convolution over addition and the assumed convolution equality give
\[
\lambda_{\mathrm L}*\Phi_\sigma
=\mu_1*\Phi_\sigma+\nu_2*\Phi_\sigma
=\nu_1*\Phi_\sigma+\mu_2*\Phi_\sigma
=\lambda_{\mathrm R}*\Phi_\sigma.
\]
% lean: Measure.add_conv

\item For any finite nonnegative Borel measure \(\lambda\) on \(\mathbb R\), define its characteristic function by
\[
\widehat{\lambda}(\xi)
:=\int_{\mathbb R}e^{i\xi\tau}\,d\lambda(\tau),
\qquad \xi\in\mathbb R.
\]
Taking characteristic functions in the preceding equality yields
\[
\widehat{\lambda_{\mathrm L}*\Phi_\sigma}(\xi)
=
\widehat{\lambda_{\mathrm R}*\Phi_\sigma}(\xi).
\]
The characteristic function of a convolution is the product of the characteristic functions, so
\[
\widehat{\lambda_{\mathrm L}}(\xi)\widehat{\Phi_\sigma}(\xi)
=
\widehat{\lambda_{\mathrm R}}(\xi)\widehat{\Phi_\sigma}(\xi)
\qquad(\xi\in\mathbb R).
\]
% lean: charFun_conv

\item The characteristic-function formula for a centered Gaussian gives
\[
\widehat{\Phi_\sigma}(\xi)
=\exp\!\left(-\frac{\sigma^2\xi^2}{2}\right)\ne0
\qquad(\xi\in\mathbb R),
\]
because the complex exponential never vanishes. This formula includes \(\sigma=0\), when the factor equals \(1\). Cancelling this nonzero factor therefore gives
\[
\widehat{\lambda_{\mathrm L}}(\xi)
=
\widehat{\lambda_{\mathrm R}}(\xi)
\qquad(\xi\in\mathbb R).
\]
% lean: charFun_gaussianReal
% lean: Complex.exp_ne_zero
% lean: mul_right_cancel₀

\item Finite Borel measures on \(\mathbb R\) are determined by their characteristic functions. Hence
\[
\lambda_{\mathrm L}=\lambda_{\mathrm R},
\qquad\text{that is,}\qquad
\mu_1+\nu_2=\nu_1+\mu_2.
\]
% lean: Measure.ext_of_charFun
\end{enumerate}
\end{proof}
